% ============================================================
%  INTRODUCTION GÉNÉRALE
% ============================================================
\chapter*{Introduction générale}
\addcontentsline{toc}{chapter}{Introduction générale}
\markboth{Introduction générale}{Introduction générale}

\vspace{0.3cm}

Ces dernières années, l'industrie aéronautique et les domaines de développement logiciel
critique connaissent une transformation profonde, poussée par des exigences croissantes
en matière de traçabilité, de qualité et de conformité aux normes internationales. La
norme \textbf{DO-178C}~\cite{do178c}, en particulier, impose aux équipes de développement un cadre
rigoureux organisant le cycle de vie du logiciel en phases distinctes : exigences de haut
niveau (HLR), exigences de bas niveau (LLR), développement (CODE), tests unitaires (LLT)
et tests d'intégration (HLT). Respecter ce processus exige une coordination précise entre
les membres de l'équipe, une traçabilité complète des tâches et un suivi rigoureux des
certifications à chaque étape.

\vspace{0.3cm}

Cependant, les outils de gestion de projets traditionnels tels que \textit{Jira}~\cite{jira},
\textit{Monday.com}~\cite{monday} ou \textit{Microsoft Project}~\cite{msproject} ne sont pas conçus pour répondre à ces
spécificités aéronautiques. Ils ne proposent pas nativement de mécanismes de gestion des
revues avec la règle Author~$\neq$~Reviewer, ni de suivi des certifications par phase, ni
d'intégration d'une intelligence artificielle capable d'analyser l'avancement en temps réel.
Par ailleurs, ils ne s'accompagnent d'aucune infrastructure DevOps~\cite{devops} permettant un déploiement
sécurisé, scalable et monitoré en production.

\vspace{0.3cm}

C'est dans ce contexte qu'a été initié le projet \textbf{SmartPM}, réalisé au sein de
\textbf{Capgemini Tunisie}~\cite{capgemini}. SmartPM est une plateforme web intelligente dédiée à la gestion
de projets aéronautiques, intégrant un assistant IA multi-modèles, un système de certifications
et revues conformes DO-178C, et une infrastructure Kubernetes permettant un déploiement
robuste et un monitoring temps réel.

\vspace{0.3cm}

Afin de mener à bien ce travail, nous avons structuré ce rapport de la manière suivante :

\begin{itemize}[leftmargin=2cm]
  \item \textbf{Le premier chapitre} présente l'organisme d'accueil Capgemini Tunisie, le
  cadre général du projet et le contexte opérationnel. Il met en évidence les limites des
  outils existants et introduit SmartPM comme réponse à ces problématiques, avant de décrire
  la méthodologie Scrum adoptée.

  \item \textbf{Le deuxième chapitre} est consacré à l'analyse des besoins et à la définition
  des exigences fonctionnelles et non fonctionnelles. Il présente également l'architecture
  technique, les choix technologiques, le backlog produit et la planification des sprints.

  \item \textbf{Le troisième chapitre} décrit la première release du projet, consacrée à la
  mise en place du socle fonctionnel de SmartPM. À travers trois sprints successifs, il détaille
  l'authentification, la gestion des projets et tâches, ainsi que les revues et certifications.

  \item \textbf{Le quatrième chapitre} traite de la deuxième release, axée sur l'intégration
  de l'intelligence artificielle et l'infrastructure DevOps. Il présente la conception et la
  réalisation des sprints 4 et 5, marquant une étape clé dans la maturité de la plateforme.
\end{itemize}

\vspace{0.3cm}

Enfin, ce rapport se conclut par une \textbf{conclusion générale} qui dresse le bilan du
travail accompli et ouvre des perspectives pour les évolutions futures de SmartPM.
